\documentclass[12pt]{article}
\usepackage{mathrsfs}
\usepackage{amsfonts} % mathbb
\usepackage{amsmath}
\usepackage{graphics}
\usepackage{latexsym}
\usepackage{amssymb}
\usepackage{color}
\usepackage{multirow}
\usepackage{amssymb}
\usepackage{bm}
\usepackage{multicol}

\newcommand{\p}{\medskip \noindent}
\newcommand{\major}[1]{\bigskip\noindent{\large{\textbf{#1}}}}
\newcommand{\minor}[1]{\medskip\noindent{\textit{\textbf{#1}}}}
\newcommand{\mathsc}[1]{\mbox{\textsc{#1}}}

\setlength{\textheight}{9in}
\oddsidemargin 0in
\evensidemargin 0in
\textwidth 6.5in
\setlength{\topmargin}{-.5in}
\newcounter{problem}
\newcommand{\pn}{\refstepcounter{problem}\theproblem}

\newcounter{subproblem}[problem]
\renewcommand{\thesubproblem}{\alph{subproblem}}
\newcommand{\spn}{\refstepcounter{subproblem}\thesubproblem}
\newcommand{\vs}{\textvisiblespace}


\begin{document}

\noindent
\textbf{CSCI 384}\\
\textbf{Regular expressions assignment}

\p
The source code for this assignment can be found in
the repository under \verb+regex+.
The file \verb+regexercises.py+ has functions for you to complete.

\p
See the programming assignment guide on Canvas to review how to 
test and turn-in your solutions.

\p
\pn. Finish the function \verb+time_format()+, which takes a string \verb+time+
and returns a boolean indicating whether \verb+time+ contains a time
formatted as \verb+hh:mm+ followed by \verb+AM+ or \verb+PM+.
Use a regular expression and the Python \verb+re+ library
(hint: \verb+re.search()+ seems to be the obvious choice here).
Use \verb+test_time_form.py+ to test;
this also serves as a complete set of examples of valid and invalid time formats.

\p
\pn. \textit{Stemming} is the task of removing the suffixes 
(and other morphemes) from a work to recover the stem or root form
of the word.
Precise stemming depends on the language and usually requires
a complicated algorithm.
We can do some crude stemming in English, however, simply by 
stripping off a common suffix like \verb+ing+, \verb+ed+, or \verb+er+.

\p
Finish the function \verb+suffix_stemmer()+ which takes a string containing a suffix
(such as \verb+ing+) and returns a string containing a regular expression to match
against a word ending with the given suffix with a capture group for the word's stem
(that is, the word excluding the suffix).
For example, if the suffix is \verb+ing+, then the regular expression returned from this function
could be used in a Python \verb+re+ library operation to match the word \verb+testing+
with \verb+test+ as the first group.

\p
Note that this function returns \textit{a string containing a regular expression} 
such as can be used in Python's regular expression library. 
Use \verb+test_suf_stem+ to test.
See that file also for examples of usage.

\p
\pn. Finish the function \verb+convert_date_format()+, which takes a string \verb+msg+
and replaces every date formatted as MM-DD-YYYY to the DD.MM.YYY format.
For example,\\
 \verb+convert_date_format("Today is 05-31-2023.")+
would return \verb+"Today is 31.5.2023."+.

\p
Use Python's \verb+re+ library. 
This will take some thinking and, possibly, research on your part. 
But the Python documentation and the examples from class and lab should help---use
those resources, not things like StackOverflow.
Hint: \verb+re.finditer()+ returns an iterator of \verb+Match+ objects.
Each \verb+Match+ object corresponds to a place where a regular expression match
occurs within a given string.
You can use \verb+Match+ methods \verb+start()+ and \verb+end()+ to find
the indices in the string where the match occurs, and you can use the \verb+Match+
object itself as a dict with the names of the match groups as keys.



\bigskip

\p
Turn in your \texttt{regexercises.py} to 
\verb+/cslab/class/cs384/[your user id]/regex+

\end{document}
